\documentclass[11pt]{article}
\usepackage[letterpaper,margin=0.85in]{geometry}\usepackage[T1]{fontenc}\usepackage[utf8]{inputenc}\usepackage{amsmath,amssymb,enumitem,booktabs,array,xcolor,fancyhdr,tabularx}
\setlist{itemsep=2pt,topsep=3pt}\pagestyle{fancy}\fancyhf{}\renewcommand{\headrulewidth}{0.3pt}
\fancyhead[L]{\small 3.1 The Chain Rule}\fancyhead[R]{\small Teacher key}\fancyfoot[C]{\small\thepage}
\definecolor{ans}{RGB}{28,125,74}\newcommand{\A}[1]{\textcolor{ans}{\textbf{#1}}}
\begin{document}

\section*{3.1 The Chain Rule --- teacher key}
\textbf{Question:} How fast does a function of a function change?\par
\subsection*{Outline and pacing (42 min)}\begin{tabular}{@{}llll@{}}\toprule Tab & Priority & Minutes & If behind\\\midrule
start & 1 & 2 & never trimmed\\
warmup & 1 & 4 & never trimmed\\
predict & 1 & 3 & never trimmed\\
explore & 3 & 8 & skipped first\\
build & 2 & 8 & shrinks to 2 of 4\\
practice & 1 & 8 & never trimmed\\
check & 1 & 4 & never trimmed\\
Huddle h1 & clock & minute 18 & pulled in by the clock\\
\bottomrule\end{tabular}
\subsection*{Targets}\begin{itemize}\item I can name the inside and the outside of a composite function.
\item I can differentiate a composite function with the chain rule.
\item I can find a slope at a point from a table.\end{itemize}
\subsection*{Predictions}\begin{itemize}
\item Expand: $y=(3x)^2=9x^2$, so the slope is $18x$. The outside rule alone gives $2(3x)=6x$. What factor is missing? \A{3, the slope of the inside.} 6x times 3 is 18x. The inside 3x changes at rate 3, and that rate multiplies in.
\item The value $y$ changes 2 units for every 1 unit of $u$. The value $u$ changes 5 units for every 1 unit of $x$. How many units does $y$ change per unit of $x$? \A{10.} The rates multiply: 2 times 5 is 10.
\end{itemize}
\section*{Huddle script (answers in green)}
\subsection*{Definitions}\begin{itemize}
\item \textbf{Composite function:} One function fed into another. The \A{inside} runs first.
\item \textbf{Chain rule:} Differentiate the \A{outside}, leave the inside, then multiply by the \A{inside?}.
\item \textbf{Reading a table:} To read h?(a), find \A{g(a)} first.
\end{itemize}
\subsection*{Worked example}
Differentiate y = sin(x\textsuperscript{2} -- 3).\\
\begin{itemize}
\item Inside: u = x\textsuperscript{2} -- 3, so du/dx = \A{2x}
\item Outside: sin u, so dy/du = \A{cos u}
\item Multiply: dy/dx = \A{2x$\cdot$cos(x\textsuperscript{2}--3)}
\end{itemize}
\subsection*{Label the assembly line}
\begin{enumerate}
\item The first machine is the \A{inside} function.
\item The last machine is the \A{outside} function.
\item The middle value u is the \A{output} of the inside.
\end{enumerate}
\section*{Every level, with answers}
\subsection*{Differentiate a composite}
\noindent\textbf{Level 1.} Differentiate. Leave the inside alone.\\ \A{Answer: $3 \left(x + 4\right)^{2}$} \quad Outside: bring down 3, lower the power. Inside $x+4$ stays. So $3(x+4)^{2}$.\par\medskip
\noindent\textbf{Level 2.} Differentiate. Multiply by the slope of the inside.\\ \A{Answer: $10 \left(5 x + 6\right)$} \quad Outside: $2(5x+6)$. Inside slope: 5. Product: $10(5x+6)$.\par\medskip
\noindent\textbf{Level 3.} Differentiate.\\ \A{Answer: $\left(-2\right) \sin\left(2 x + 5\right)$} \quad Outside: $-\sin(2x+5)$. Inside slope: 2. Product: $-2\sin(2x+5)$.\par\medskip
\noindent\textbf{Level 4.} Differentiate. Rewrite the root as a power if that helps.\\ \A{Answer: $\frac{x}{\sqrt{x^{2} + 2}}$} \quad $\dfrac{1}{2\sqrt{x^{2}+2}}\cdot 2x=\dfrac{x}{\sqrt{x^{2}+2}}$.\par\medskip
\subsection*{Name the inside}
\noindent\textbf{Level 1.} In $y=(x+4)^{3}$, the inside function is...\\ \A{(A) x + 4} \quad (B) x\textsuperscript{3} \quad (C) y \quad (D) 4\\ \A{Answer: A} \quad x goes into $x+4$ first; then the result is raised to the power 3.\par\medskip
\noindent\textbf{Level 2.} In $y=\sin(2x)$, which function is the outside?\\ (A) 2x \quad \A{(B) sin} \quad (C) 2 \quad (D) x\\ \A{Answer: B} \quad x is multiplied by 2 first. Sine is applied last, so sine is the outside.\par\medskip
\noindent\textbf{Level 3.} In $y=e^{\sin x}$, the inside function is...\\ (A) e? \quad (B) x \quad (C) e\textasciicircum{}(sin x) \quad \A{(D) sin x}\\ \A{Answer: D} \quad x goes into sine first. The exponential is applied to that output.\par\medskip
\noindent\textbf{Level 4.} How many factors does the derivative of $y=\sin\left(\sqrt{x^{2}+6}\right)$ have?\\ (A) Two \quad (B) One \quad \A{(C) Three} \quad (D) Four\\ \A{Answer: C} \quad Three layers: sine, the square root, and $x^{2}+6$. Each gives one factor.\par\medskip
\subsection*{Read a table}
\noindent\textbf{Level 1.} A table gives $g(2)=3$, $g'(2)=4$ and $f'(3)=5$. Let $h(x)=f(g(x))$. Find $h'(2)$.\\ \A{Answer: 20} \quad $h'(2)=f'(g(2))\cdot g'(2)=f'(3)\cdot g'(2)=5\cdot4=20$.\par\medskip
\noindent\textbf{Level 2.} A table gives $f'(1)=3$, $g(1)=6$, $g'(1)=-5$ and $f'(6)=2$. Let $h(x)=f(g(x))$. Find $h'(1)$.\\ \A{Answer: -10} \quad g(1) = 6, so $h'(1)=f'(6)\cdot g'(1)=2\cdot(-5)=-10$. The value $f'(1)$ is a distractor.\par\medskip
\noindent\textbf{Level 3.} A table gives $f(4)=1$, $f'(4)=4$, $g'(4)=6$ and $g'(1)=2$. Let $h(x)=g(f(x))$. Find $h'(4)$.\\ \A{Answer: 8} \quad f(4) = 1, so $h'(4)=g'(1)\cdot f'(4)=2\cdot(4)=8$.\par\medskip
\noindent\textbf{Level 4.} A table gives $g(4)=6$, $g'(4)=2$, $f(6)=4$ and $f'(6)=5$. Let $h(x)=f(g(x))$. The tangent line to $h$ at $x=4$ is used to estimate $h(5)$. Find the estimate.\\ \A{Answer: 14} \quad $h(4)=f(6)=4$ and $h'(4)=5\cdot(2)=10$. One step right gives $4+(10)=14$.\par\medskip
\subsection*{Slope at a point}
\noindent\textbf{Level 1.} Let $y=(x+2)^{2}$. Find the slope of the curve at $x=2$.\\ \A{Answer: 8} \quad $y'=2(x+2)$, so at $x=2$ the slope is $2\cdot4=8$.\par\medskip
\noindent\textbf{Level 2.} Let $y=(4x+3)^{2}$. Find the slope of the curve at $x=1$.\\ \A{Answer: 56} \quad $y'=8(4x+3)$, so at $x=1$ the slope is $8\cdot7=56$.\par\medskip
\noindent\textbf{Level 3.} Let $y=\ln(x^{2}+8)$. Find the slope of the curve at $x=1$.\\ \A{Answer: 0.2222} \quad $y'=\dfrac{2x}{x^{2}+8}$, so at $x=1$ the slope is $\dfrac{2}{9}$.\par\medskip
\noindent\textbf{Level 4.} Let $y=\sqrt{x^{2}+33}$. Find the slope of the curve at $x=4$.\\ \A{Answer: 0.5714} \quad $y'=\dfrac{x}{\sqrt{x^{2}+33}}$, so at $x=4$ the slope is $\dfrac{4}{\sqrt{49}}=0.5714$.\par\medskip
\section*{Common misconceptions}\begin{tabularx}{\textwidth}{@{}X X@{}}\toprule Students say & Say instead\\\midrule
The derivative of sin(x\textsuperscript{2}) is cos(x\textsuperscript{2}). & The inside also changes. Multiply by its slope, 2x.\\
The derivative of (3x + 1)? is 4(3x + 1)\textsuperscript{3}. & That is the outside only. Multiply by the inside slope, 3.\\
Read h?(a) as f?(a) $\cdot$ g?(a). & f? is read at g(a), the output of the inside. Find g(a) first.\\
Stop after one layer of cos\textsuperscript{3}(4x). & Peel every layer: the power, then cosine, then 4x. Each gives a factor.\\
\bottomrule\end{tabularx}
\end{document}